
Om rollen
Vill du arbeta i ett teknikdrivet bolag där utveckling, data och prestanda står i centrum? Car.info söker nu en DevOps Engineer som vill vara med och vidareutveckla en modern SaaS-plattform i en miljö där teknikval och förbättringsarbete drivs nära utvecklingsteamet.

I rollen arbetar du nära utvecklare och produktteam för att säkerställa stabila, skalbara och effektiva lösningar. Du blir en viktig del i att förbättra utvecklingsflöden, automatisera processer och skapa en robust infrastruktur där systemen är snabba, tillgängliga och driftsäkra.

Det är viktigt att förtydliga att detta inte är en renodlad DevOps-roll. Rollen är bred och innehåller flera tekniska ansvarsområden där DevOps är en del av helheten. Arbetet innefattar även utveckling och tekniskt förbättringsarbete, där PHP förekommer som en naturlig del av vardagen. Rollen passar därför dig som trivs i en varierad teknisk roll och uppskattar att arbeta brett, snarare än att enbart fokusera på DevOps.
Kvalifikationer
Erfarenhet av DevOps eller liknande teknisk roll
Erfarenhet av CI/CD och automatiserade deployprocesser
Erfarenhet av molnplattformar (AWS, Azure eller GCP)
Erfarenhet av containerteknologier som Docker och Kubernetes
Erfarenhet av scripting eller programmering (exempelvis Bash eller Python)
Erfarenhet av PHP är ett krav
God förståelse för systemarkitektur, prestanda och skalbarhet
Flytande svenska och engelska i tal och skrift
Meriterande
Erfarenhet av övervakning och loggning (Prometheus, Grafana, ELK)
Infrastructure as Code
Erfarenhet från SaaS- eller produktbolag
Arbetsuppgifter
Designa, implementera och underhålla CI/CD-pipelines
Automatisera utvecklings-, test- och distributionsprocesser
Hantera och optimera molninfrastruktur
Bidra till utveckling och underhåll av system där PHP ingår
Säkerställa hög tillgänglighet, prestanda och säkerhet i systemen
Samarbeta nära utvecklingsteam för effektiva leveranser
Identifiera och driva tekniska förbättringar i plattformen
Företaget erbjuder
Car.info är ett teknikdrivet och snabbt växande bolag med bas i Ängelholm. De utvecklar smarta SaaS-lösningar med fokus på användarvänlighet, prestanda och kundnytta. Här kombineras ett modernt arbetssätt med hög teknisk kompetens och stort affärsdriv.
Central roll i ett växande företag med höga ambitioner
Möjlighet att påverka både teknik och produktutveckling
Delaktighet i tekniska beslut och förbättringsarbete
Kort väg från idé till beslut – snabba och synliga resultat
Direkt feedback från användare och möjlighet att påverka produkten i realtid
En utvecklarvänlig arbetsmiljö med fokus på innovation, kvalitet och samarbete
Förmåner
Marknadsmässig lön (med möjlighet för årlig bonus).
5 veckors semester (6 veckor för dig som fyllt 40) med ytterligare intjänad semesterdag vartannat år.
Centralt kontor i vackra Ängelholm.
Möjlighet att arbeta hemifrån 2 dagar i veckan.
Ett härligt team av kollegor.
Gratis kaffe/frukt/läsk, måndagsfrukost och fredagsfika.
Ansökan


Är du redo att bli en del av vårt team?
Skicka ditt CV och personliga brev till:
📩 JC@happygroup.se

Frågor om tjänsten besvaras av:
📧 melvin@happygroup.se
Urval sker löpande – vänta inte med att skicka in din ansökan.